\documentclass[11pt,letterpaper]{article}
\usepackage[T1]{fontenc}
\usepackage[utf8]{inputenc}
\usepackage{lmodern,amsmath,amssymb,amsthm,mathtools,mathrsfs}
\usepackage[margin=1in]{geometry}
\usepackage{microtype,booktabs,longtable,array,enumitem,xcolor,xurl}
\usepackage{fancyhdr}
\usepackage[colorlinks=true,linkcolor=blue!40!black,citecolor=blue!40!black,urlcolor=blue!40!black]{hyperref}
\usepackage{bookmark}
\setlength{\headheight}{14pt}
\setlength{\emergencystretch}{3em}
\pagestyle{fancy}\fancyhf{}
\fancyhead[L]{\small Finite polylogarithm alphabets for Hankel sums}
\fancyhead[R]{\small ProveIt research continuation}
\fancyfoot[C]{\thepage}
\numberwithin{equation}{section}
\newtheorem{theorem}{Theorem}[section]
\newtheorem{lemma}[theorem]{Lemma}
\newtheorem{proposition}[theorem]{Proposition}
\newtheorem{corollary}[theorem]{Corollary}
\theoremstyle{definition}\newtheorem{definition}[theorem]{Definition}
\theoremstyle{remark}\newtheorem{remark}[theorem]{Remark}
\newtheorem{example}[theorem]{Example}
\newcommand{\Q}{\mathbb Q}\newcommand{\C}{\mathbb C}\newcommand{\N}{\mathbb N}\newcommand{\Z}{\mathbb Z}
\newcommand{\dd}{\,\mathrm d}\newcommand{\Li}{\operatorname{Li}}
\newcommand{\wt}{\operatorname{wt}}\newcommand{\Res}{\operatorname{Res}}
\newcommand{\HH}{\mathcal H}\newcommand{\FP}{\operatorname{FP}}
\newcommand{\tht}{\theta}\newcommand{\fall}[2]{#1^{\underline{#2}}}
\newcommand{\src}[1]{\nolinkurl{#1}}
\newcommand{\snapshot}{505c62fddea6b4a9eb61fee0c76e6c021e9c7128}
\hypersetup{pdftitle={A Finite Polylogarithm Alphabet for Rational Hankel Sums},pdfauthor={Research prepared with ChatGPT for Vladimir Reshetnikov},pdfsubject={Confluent poles, harmonic determinants, Gamma coefficients, multiple polylogarithms, and Stieltjes spectral identities}}
\title{\textbf{A Finite Polylogarithm Alphabet\\for Rational Hankel Sums}\\[0.45em]
\large Confluent pole coefficients, Gamma-product normalizations,\\endpoint cancellation, and Stieltjes spectral identities}
\author{Research prepared for Vladimir Reshetnikov\\\small Mathematical derivations and computational checks with ChatGPT}
\date{11 October 2026}
\begin{document}
\maketitle
\begin{abstract}
We give an explicit finite multiple-polylogarithm reduction for every
cutoff Hankel determinant of a rational weight whose poles are at
nonpositive integers. The number of particles may be arbitrary. The
controlling invariant is not determinant size or ordinary harmonic
rank: a pole of multiplicity $s$ supplies the decreasing positive
sequence $s,s-2,s-4,\ldots$. The sum of the $r$ largest available entries
bounds the polylogarithmic weight for an $r$-particle determinant.
A confluent Vandermonde factorization proves this statement; a
Hermite-reduction argument turns it into a terminating identity compiler.
For a single pole the formal bound is attained at every rank, and a
specified extremal coefficient is an explicit product of Gamma factors.
For several simple poles we derive a fixed-size Hahn determinant for
the highest harmonic coefficient, a complete finite formula for each
endpoint principal part, and an exact spectral description of one-residue
exceptional parameters. A separate spectral determinant identity yields
ordinary all-order relations involving generalized Stieltjes constants
and derivatives of depth-two and higher zeta functions. All infinite
identities are proved analytically; exact finite tests and numerical
diagnostics are supplied separately. Historical priority beyond the
inspected sources is not claimed.
\end{abstract}
\medskip
\noindent\textbf{Main distinction.} ``Weight'' below is a grading of an
explicit expression, not a claim of transcendence, algebraic independence,
or minimal multiple-polylogarithm depth. The unrestricted harmonic-tail
moment-prefix conjecture and the manuscript's Gaussian $S_6$ and revised
$S_8$ identities are not resolved here.
\clearpage
\tableofcontents
\newpage

\section{Research scope and comparison with the repository}

The natural starting point for this continuation is the several-pole
question in the incoming gap-six report, rather than an older Gamma
normal-form question that has already been answered. In the inspected
\emph{Gap Six Harmonic Determinants} package, Section~3 gives an
arbitrary-particle formula for two adjacent simple poles using only a
dilogarithm and powers of $-\log(1-z)$. Section~5, Questions~Q7--Q9,
asks for several-pole extensions, a classification of when classical
polylogarithms suffice, and exact Euler primitives \cite{gap}.
The same package already proves the complex rational-weight
\emph{ordinary harmonic degree--rank} theorem. We do not claim that
result again.

The new contribution developed here is a different grading and its exact
analytic realization. Repeated poles make ordinary harmonic degree
insufficient: $H_M^{(1)}H_M^{(3)}-(H_M^{(2)})^2$ has harmonic degree two
but weight four. Our confluent minor calculation gives a uniform finite
alphabet, explicitly tracks the cutoff poles needed by the compiler,
and supplies a nonzero extremal coefficient for the single-pole family.
The several-simple-pole result gives more than a bare appeal to a
Christoffel transformation: it specifies the compressed measure, the
normalization, the principal-part cancellation equations, and a
positive-matrix description of exceptional residue parameters.

Several ingredients are classical. Finite Cauchy--Binet, confluent
Vandermonde matrices, finite harmonic quasi-shuffles \cite{hoffman},
Christoffel determinant transformations \cite{krattenthaler}, uniform
Hahn polynomials \cite{dlmf-hahn}, and integration of hyperlogarithms
against rational functions \cite{panzer} are used with their original
status intact. Determinantal stability is also an established method
\cite{borcea}. The contribution is their particular combination with
the pole-multiplicity invariant, its sharp coefficient, and the
resulting identity families. Neither the software architecture for
hyperlogarithmic integration nor the classical weight-five Euler sums
used in an example are claimed as new.

\subsection{A fixed comparison point and a moving repository}

During retrieval, \texttt{main} changed while older incoming archives
were being placed in report directories. The final observation used here
is the exact commit
\begin{center}\small\texttt{\snapshot}.\end{center}
The gap-six archive at that commit has Git blob
\src{d7695b94c3117d7bdcb5e75fcf930c410e9c31b6}.
Its sections on complex rank, adjacent-pole generators, and further
questions were read directly from the archive. The universal Gamma
normal-form and Gamma-normalization archives were also checked to avoid
reopening their completed questions \cite{universal,normalizations}.
The source audit accompanying this article identifies the scope of each
inspection. It is not an assertion that every page in the repository
or every incoming archive was reviewed. No repository file was modified.

\section{Conventions and the exact object to be reduced}

Let $r\ge1$ and $M\ge0$ be integers. We use the nodes $1,\ldots,M$ and
write
\begin{equation}\label{eq:R}
 R(x)=P(x)+\sum_{a\in A}\sum_{\ell=1}^{s_a}
                     \frac{c_{a,\ell}}{(x+a)^\ell},
 \qquad A\subset\Z_{\ge0}\text{ finite}.
\end{equation}
Here $P$ is a polynomial, $c_{a,s_a}\ne0$, and coefficients may be
complex. All denominators are therefore nonzero on the summation nodes.
The field $K$ is generated over $\Q$ by the coefficients of $R$.
Set
\begin{align}
 \mu_j(M)&=\sum_{n=1}^M n^j R(n),\notag\\
 D_{r,M}(R)&=\det(\mu_{i+j}(M))_{0\le i,j<r},
 &G_r(z;R)&=\sum_{M\ge1}D_{r,M}(R)z^M.
 \label{eq:det}
\end{align}
The empty-node matrix gives $D_{r,0}=0$. In fact $D_{r,M}=0$ when
$M<r$. Finite Cauchy--Binet gives the exact identity
\begin{equation}\label{eq:ensemble}
 D_{r,M}(R)=\sum_{1\le n_1<\cdots<n_r\le M}
             \prod_{i<j}(n_j-n_i)^2\prod_{j=1}^r R(n_j).
\end{equation}
The source's zero-based nodes $0,\ldots,N$ translate to these conventions
by $M=N+1$, $n_{\mathrm{source}}=n-1$, and
$a_{\mathrm{source}}=a+1$. A monic affine change of polynomial basis has
determinant one, so it preserves the Hankel determinant.
There is no factor $r!$ in this ordered sum. Expansion of the two
Vandermonde determinants proves \eqref{eq:ensemble} directly.
The terms grow at most polynomially in $M$, so \eqref{eq:det} converges
absolutely for $|z|<1$, locally uniformly there.

The generalized harmonic sums are
\begin{equation}\label{eq:h}
 H_M^{(p)}=\sum_{n=1}^M n^{-p},\qquad
 h_{a,p}(M)=\sum_{n=1}^M(n+a)^{-p}.
\end{equation}
A polynomial expression in formal symbols $h_{a,p}$ has
\emph{weighted harmonic degree} obtained by assigning
$\wt(h_{a,p})=p$ and weight zero to coefficients depending polynomially
on $M$. This grading is first imposed before identifying the shifted
harmonic symbols.
For integral $a\ge0$, the exact shift rule is
\begin{equation}\label{eq:shift}
 h_{a,p}(M)=H_M^{(p)}+
           \sum_{b=1}^a\frac1{(M+b)^p}-H_a^{(p)}.
\end{equation}
Thus an inverse cutoff factor of order $p$ inherits a cost of $p$.
The constants $H_a^{(p)}$ are rational and have cost zero.

We use the decreasing-index convention
\begin{equation}\label{eq:mpl}
 \Li_{p_1,\ldots,p_k}(z)=
 \sum_{n_1>\cdots>n_k\ge1}
       \frac{z^{n_1}}{n_1^{p_1}\cdots n_k^{p_k}},
 \qquad p_i\in\Z_{>0}.
\end{equation}
Its weight is $p_1+\cdots+p_k$. Denote by $\mathscr L_W$ the span over
$K(z)$ of $1$ and all these functions of weight at most $W$. Every
rational prefactor constructed below has poles only at $0$ and $1$.
No assertion that the displayed generators are a minimal basis is needed.

\section{The confluent pole-multiplicity theorem}

\subsection{A moment factorization by local derivative columns}

For a fixed pole $-a$ of order $s=s_a$, define the $r\times s$ matrix
\begin{equation}\label{eq:V}
 V_a(i,p)=\begin{cases}\binom ip(-a)^{i-p},&p\le i,\\0,&p>i,
                    \end{cases}
 \qquad 0\le i<r,\quad0\le p<s.
\end{equation}
The value $0^0$ occurring here is one. Define a symmetric $s\times s$
harmonic matrix by
\begin{equation}\label{eq:C}
 C_a(p,q)=\sum_{\ell=p+q+1}^{s}
                   c_{a,\ell}h_{a,\ell-p-q}.
\end{equation}
An empty sum is zero. In particular
$\wt C_a(p,q)\le s-p-q$ for a nonzero entry.

\begin{lemma}[Confluent harmonic factorization]\label{thm:factor}
There is a matrix $B_r(M)$ with entries in $K[M]$ such that the formal
moment matrix is
\begin{equation}\label{eq:factor}
       \boldsymbol\mu_r(M)=B_r(M)+\sum_{a\in A}V_a C_a V_a^{\mathsf T}.
\end{equation}
The matrix $B_r$ is the finite sum of the polynomial parts of
$x^{i+j}R(x)$. The transpose is ordinary transpose, not conjugate
transpose.
\end{lemma}
\begin{proof}
Expand $x^i$ and $x^j$ at $x=-a$:
\[
 x^i=\sum_p\binom ip(-a)^{i-p}(x+a)^p.
\]
After multiplying and dividing by $(x+a)^\ell$, terms with $p+q<\ell$
sum to $h_{a,\ell-p-q}$ and are exactly \eqref{eq:C}--\eqref{eq:factor}.
The other terms are nonnegative powers of $x+a$ and have polynomial
Faulhaber sums in $M$. Equivalently, the partial-fraction coefficient
uses
$\binom{i+j}{t}=\sum_{p+q=t}\binom ip\binom jq$.
This proves every matrix entry, including the polynomial part of $R$.
\end{proof}

\subsection{The invariant and its closed combinatorial form}

For $r\ge0$, define
\begin{equation}\label{eq:capacity}
 W_r(\mathbf s)=
 \max_{\substack{0\le k_a\le s_a\\\sum_a k_a\le r}}
                 \sum_{a\in A}k_a(s_a-k_a+1).
\end{equation}
For no poles, $W_r=0$.

\begin{theorem}[Weighted harmonic degree]\label{thm:weight}
Every monomial in the formal determinant obtained from
\eqref{eq:factor} has weighted harmonic degree at most
$W_r(\mathbf s)$. Equivalently, form the multiset
\begin{equation}\label{eq:marginals}
 \bigsqcup_{a\in A}\{s_a,s_a-2,s_a-4,\ldots\}\cap\Z_{>0},
\end{equation}
retaining multiplicities. Then $W_r$ is the sum of the largest $r$
entries, or all entries if there are fewer than $r$. In particular
\begin{equation}\label{eq:saturation}
 W_r\le W_\infty:=\sum_{a\in A}
                           \left\lfloor\frac{(s_a+1)^2}{4}\right\rfloor.
\end{equation}
These are bounds on an exact expression grading, not inequalities
between values of special functions.
\end{theorem}
\begin{proof}
Concatenate the $V_a$ and make the $C_a$ into a block diagonal matrix
$C$. In the expansion of $\det(B+VCV^{\mathsf T})$, choose $k$ columns
from the second summand and the complementary columns from $B$.
Apply Cauchy--Binet to the resulting $k\times k$ minors of
$VCV^{\mathsf T}$. All dependence on harmonic symbols lies in minors
$\det C[I,J]$ with $|I|=|J|=k\le r$.

A nonzero minor of this block diagonal matrix selects the same number
$k_a$ of row and column indices from every pole block. Let $I_a,J_a$
be those local indices, each a subset of $\{0,\ldots,s_a-1\}$.
Every term of its determinant has weight at most
\[
 k_a s_a-\sum_{p\in I_a}p-\sum_{q\in J_a}q
 \le k_a s_a-k_a(k_a-1)=k_a(s_a-k_a+1).
\]
The second inequality uses the smallest possible sum of $k_a$ distinct
nonnegative integers, $k_a(k_a-1)/2$, separately for rows and columns.
The $V$ and $B$ minors have weight zero. Adding the block contributions
and maximizing proves \eqref{eq:capacity}.

The increase in $k(s-k+1)$ when $k$ is replaced by $k+1$ is $s-2k$.
These marginal increases decrease within each pole block. Maximizing
with at most $r$ choices therefore selects the largest positive
marginals globally; their decreasing order automatically respects every
block's prefix constraint. Zero marginals may be omitted and negative
ones never help. Summing all positive marginals yields
$\lfloor(s+1)^2/4\rfloor$ for each block. This proves
\eqref{eq:marginals} and \eqref{eq:saturation}.
\end{proof}

For simple poles this specializes to $W_r=\min(r,|A|)$.
For one pole of orders $1,2,3,4,5,6,7$, the stabilized weights are
respectively $1,2,4,6,9,12,16$. The jump from order two to order three
is real: pole multiplicity itself is not the correct weight budget.

\begin{remark}[What cancellation can and cannot imply]
The upper bound permits special residues to cancel its highest
components. It does not assert weighted-degree equality for every
multiple-pole complex weight. It also does not prove that a function
of weight $W_r$ cannot be written using functions of smaller weight.
A precise formal sharpness theorem for a single pole follows next.
\end{remark}

\section{A sharp single-pole coefficient with Gamma normalization}

Put $R(x)=x^{-s}$ with $s\ge1$ integral. Polynomial division is trivial:
\begin{equation}\label{eq:singlemom}
 \mu_j(M)=\begin{cases}H_M^{(s-j)},&j<s,\\
 S_{j-s}(M),&j\ge s,\end{cases}
 \qquad S_t(M)=\sum_{n=1}^M n^t.
\end{equation}
Let $\mathcal P_{r,s}(M;H_1,\ldots,H_s)$ be this formal determinant.
Write
\begin{equation}\label{eq:uniformnorm}
 \mathcal U_q(M)=\prod_{j=0}^{q-1}
       \frac{\Gamma(M+j+1)}
       {(2j+1)\Gamma(M-j)\binom{2j}{j}^{\!2}},
 \qquad \mathcal U_0(M)=1.
\end{equation}
For $M\ge q$ all factors are finite and positive. They also have
polynomial continuations because each Gamma quotient is
$(M-j)_{2j+1}$.

\begin{theorem}[An explicit nonzero extremal coefficient]\label{thm:sharp}
Let
\[
 k=\min\bigl(r,\lceil s/2\rceil\bigr),\qquad q=r-k,
 \qquad W=k(s-k+1)=W_r(s).
\]
The coefficient of $H_{s-k+1}^{\,k}$ in $\mathcal P_{r,s}$ is
\begin{equation}\label{eq:sharp}
 \boxed{\displaystyle
 [H_{s-k+1}^{\,k}]\mathcal P_{r,s}=
 (-1)^{k(k-1)/2}\mathcal U_q(M)
 \begin{cases}
  1,&q=0\text{ or }s\text{ even},\\[2pt]
  \dfrac{(M+1)_q}{\binom{2q}{q}},&q>0\text{ and }s\text{ odd}.
 \end{cases}}
\end{equation}
Consequently the formal weighted harmonic degree is exactly $W_r(s)$
as a polynomial in $M$, and also at every integer $M\ge r$.
\end{theorem}
\begin{proof}
A term contributing the displayed monomial must select $k$ harmonic
entries. Its harmonic weight is $k(s-k+1)$. Equality in the minor bound
in Theorem~\ref{thm:weight} forces the chosen row and column sets to be
$\{0,\ldots,k-1\}$. Within that harmonic minor, the monomial with all
indices $s-k+1$ requires $i+\sigma(i)=k-1$ at every row. The unique
permutation is reversal and has sign $(-1)^{k(k-1)/2}$.

The complementary minor has indices $i,j\ge k$. If it is nonempty,
then $k=\lceil s/2\rceil$ and $2k-s=\varepsilon\in\{0,1\}$.
It is the polynomial moment determinant
\begin{equation}\label{eq:sharpgram}
       \det\left(\sum_{n=1}^M n^{i+j+\varepsilon}\right)_{0\le i,j<q}.
\end{equation}
For $\varepsilon=0$, its value is the product of the monic uniform
Hahn norms, namely $\mathcal U_q(M)$.
For $\varepsilon=1$, the one-factor Christoffel formula gives
$(-1)^q\mathcal U_q(M)p_q(0;M)$, where $p_q$ is monic and orthogonal
on $1,\ldots,M$ with unit weights. Its endpoint value is
\[
       p_q(0;M)=(-1)^q\frac{(M+1)_q}{\binom{2q}{q}}.
\]
The finite Hahn formula and normalization are recalled in
Section~\ref{sec:hahn}; they are classical. This proves
\eqref{eq:sharp}. Positivity of \eqref{eq:sharpgram} when $M\ge q$
proves nonvanishing in the stated range. The upper bound already proved
then gives exact formal degree. When $s$ is even, an additional
zero-marginal choice can have the same weight and a different harmonic
degree; it cannot contribute the monomial of degree $k$ selected here.
\end{proof}

For example, the coefficients of $H_2^2$ for the triple pole are
\begin{align}
 [H_2^2]\mathcal P_{2,3}&=-1,\notag\\
 [H_2^2]\mathcal P_{3,3}&=-\frac{M(M+1)}2,\notag\\
 [H_2^2]\mathcal P_{4,3}&=
 -\frac{M^2(M-1)(M+1)^2(M+2)}{72}.
 \label{eq:sharpexamples}
\end{align}
All ranks nevertheless have the same maximal weight four. This is a
coefficient identity with an explicit Gamma-product source, not an
asymptotic observation.

\section{From the pole grading to a finite polylogarithm identity}
\label{sec:compiler}

\subsection{Cost-preserving cutoff partial fractions}

\begin{lemma}[Normal form after merging integer shifts]\label{thm:cutoffnormal}
After applying \eqref{eq:shift}, $D_{r,M}(R)$ is a finite sum of terms
of the following two types:
\begin{equation}\label{eq:cutoffterms}
       p(M)\prod_{j=1}^v H_M^{(b_j)},\qquad
       \frac{c}{(M+h)^e}\prod_{j=1}^v H_M^{(b_j)},
\end{equation}
where $p\in K[M]$, $c\in K$, $h\in\{1,\ldots,\max A\}$,
$e\ge1$, and
\begin{equation}\label{eq:cost}
 \sum_j b_j\le W_r\quad\text{in the first type},\qquad
 e+\sum_jb_j\le W_r\quad\text{in the second type}.
\end{equation}
The empty harmonic product is one. If $A$ is empty or $A=\{0\}$,
there are no cutoff-denominator terms.
\end{lemma}
\begin{proof}
Before merging shifts, Theorem~\ref{thm:weight} supplies a polynomial
in the $h_{a,p}$ with polynomial coefficients in $M$. Expanding each
factor by \eqref{eq:shift}, choose either a harmonic factor of weight
$p$, an endpoint denominator of order $p$, or a rational constant.
The total of retained harmonic weight and all denominator orders cannot
increase. A rational coefficient with total denominator order $E$ has
only the listed poles, each of order at most $E$. Partial fractions
and polynomial division give \eqref{eq:cutoffterms}; every resulting
individual pole order is at most $E$. This proves \eqref{eq:cost}.
No cancellation assumption is used.
\end{proof}

\subsection{Finite harmonic products and their generating functions}

For a word $\mathbf b=(b_1,\ldots,b_v)$ of positive integers define
\[
 H_M(\mathbf b)=\sum_{M\ge n_1>\cdots>n_v\ge1}
                        n_1^{-b_1}\cdots n_v^{-b_v}.
\]
The empty word has value one. Partitioning a product of finite sums by
which indices coincide gives
\begin{equation}\label{eq:stuffle}
 H_M(a)H_M(b)=H_M(a,b)+H_M(b,a)+H_M(a+b),
\end{equation}
and the recursive rule
\begin{equation}\label{eq:stufflerec}
 (a,u)*(b,v)=(a,u*(b,v))+(b,(a,u)*v)+(a+b,u*v).
\end{equation}
The empty word is the identity and identical output words retain their
integer multiplicities. This is the positive finite-sum quasi-shuffle,
not the version with a negative merger sign used for some star-sum
conventions \cite{hoffman}. The rule preserves total weight.
Furthermore, for every nonempty word,
\begin{equation}\label{eq:hgen}
 \sum_{M\ge1}H_M(\mathbf b)z^M
               =\frac{\Li_{\mathbf b}(z)}{1-z}.
\end{equation}
Indeed the largest summation index $n_1$ is contained in every cutoff
$M\ge n_1$; summing its geometric tail proves the formula. For the empty
word the corresponding generator is $z/(1-z)$, not $1/(1-z)$.
Absolute convergence for $|z|<1$ justifies these rearrangements.

\subsection{A closed rational integration step}

Use hyperlogarithm words on the two letters $0,1$, with kernels
$f_0(z)=1/z$ and $f_1(z)=1/(1-z)$. Let $I_\emptyset=1$ and, for a
nonempty word ending in $1$, set
\[
 I_{a\omega}(z)=\int_0^z f_a(t)I_\omega(t)\dd t.
\]
These anchored integrals are equivalent to \eqref{eq:mpl}:
\begin{equation}\label{eq:words}
 \Li_{b_1,\ldots,b_v}(z)=
 I_{0^{b_1-1}1\cdots0^{b_v-1}1}(z).
\end{equation}
Their weight is word length. Intermediate rational integrations can
produce $\log z$; the following statement includes the necessary
cancellation and anchoring convention.

\begin{lemma}[Rational integration without a new letter]\label{thm:primitive}
Let $F$ be a finite sum of rational multiples of words ending in $1$
and of the empty word. Suppose every rational factor has poles only
at $0,1$ and all word lengths are at most $w$. If $F$ is holomorphic
at zero, its primitive vanishing at zero is a rational linear
combination of words ending in $1$ of lengths at most $w+1$ and the
empty word. It can be constructed with exact rational operations and
finite differentiation. The same holds for $F(z)/z$ if $F(0)=0$.
\end{lemma}
\begin{proof}
For a rational function with poles only at $0,1$, ordinary partial
fractions give the exact Hermite decomposition
\begin{equation}\label{eq:hermite}
 R(z)=Q'(z)+\frac{\alpha}{z}+\frac{\beta}{1-z},
 \qquad Q\in K(z),\quad\alpha,\beta\in K.
\end{equation}
Polynomial terms and poles of order at least two integrate into $Q$.
For a nonempty word $\omega=a\nu$, integration by parts gives
\begin{equation}\label{eq:intstep}
 \int R I_\omega\dd z
 =Q I_\omega+\alpha I_{0\omega}+\beta I_{1\omega}
             -\int Q f_a I_\nu\dd z.
\end{equation}
The remaining integral has a strictly shorter word and a rational
factor with the same possible poles. Repeated application terminates.
The two prepended-letter terms have length at most $w+1$; every other
term has no greater length. For the empty word the primitive is
$Q+\alpha\log z-\beta\log(1-z)$.

Combine equal words and simplify rational coefficients before taking
an endpoint. The constructed primitive is a meromorphic germ plus a
constant multiple of $\log z$, since all other retained words are
analytic at zero. Its derivative is the assumed holomorphic $F$.
Therefore its Laurent principal part and its $\log z$ coefficient
vanish. The remaining constant is determined by finitely many Taylor
coefficients: each rational factor has a finite pole order. Subtract
that constant. This gives the unique anchored analytic primitive and
proves the claim. The case $F/z$ follows because it is holomorphic
under the additional hypothesis. This is the two-letter specialization
of the standard hyperlogarithmic integration method \cite{panzer}.
\end{proof}

For $\theta=z\,d/dz$ and an integer $h\ge0$, define
\begin{equation}\label{eq:inverse}
 \mathcal J_hF(z)=z^{-h}\int_0^z t^{h-1}F(t)\dd t.
\end{equation}
If $h=0$, require $F(0)=0$. If $h>0$ and $F$ is analytic, the
apparent singularity is removable. The coefficient identity is
\begin{equation}\label{eq:inversecoeff}
 \mathcal J_h\!\left(\sum_{M\ge1}a_Mz^M\right)
              =\sum_{M\ge1}\frac{a_M}{M+h}z^M.
\end{equation}
Differentiation by $\theta$ does not increase word length; one
application of $\mathcal J_h$ increases it by at most one. There is no
unspecified integration constant in \eqref{eq:inverse}.

\begin{theorem}[Uniform finite polylogarithm alphabet]\label{thm:compiler}
For every rational weight \eqref{eq:R} and every $r\ge1$,
\begin{equation}\label{eq:maincompiler}
 \boxed{\displaystyle G_r(z;R)\in\mathscr L_{W_r(\mathbf s)},
                                    \qquad |z|<1.}
\end{equation}
An explicit terminating construction is as follows. Expand the
confluent determinant, merge shifts, put its coefficients in the
partial-fraction form \eqref{eq:cutoffterms}, and quasi-shuffle each
harmonic product. Replace it by \eqref{eq:hgen}. Replace multiplication
by $p(M)$ by $p(\theta)$, and division by $(M+h)^e$ by
$\mathcal J_h^e$. Apply \eqref{eq:hermite}--\eqref{eq:intstep} to remove
all unevaluated integrals. The output has rational coefficients in
$K(z)$ and no letters other than $0,1$.
\end{theorem}
\begin{proof}
A harmonic product of weight $b$ gives a finite sum of polylogarithms
of weight $b$ divided by $1-z$. The empty product has the rational
generator already specified. Polynomial coefficient operators do not
increase weight. Each inverse-cutoff factor raises it by at most its
order, by Lemma~\ref{thm:primitive}. Lemma~\ref{thm:cutoffnormal}
therefore proves the stated weight bound term by term. All operations
are finite: the determinant has finite size, partial fractions are
finite, the quasi-shuffle has finitely many words, and the integration
recursion strictly decreases word length. Identities
\eqref{eq:hgen} and \eqref{eq:inversecoeff} prove equality of all
Taylor coefficients. Both sides are analytic in $|z|<1$, so the
constructed expression equals the original generator throughout that
disk, with removable values at zero taken collectively.
\end{proof}

\begin{corollary}[Every quadratic denominator uses only a dilogarithm]
\label{thm:quadratic}
Suppose the reduced denominator of $R$ has degree at most two and
all its roots are nonpositive integers; the polynomial part is arbitrary.
For every $r\ge1$ there are rational functions $A_r,B_r,C_r,E_r\in K(z)$
with poles only at $0,1$ such that, with $L=-\log(1-z)$,
\begin{equation}\label{eq:quadratic}
       G_r(z;R)=A_r(z)+B_r(z)L+C_r(z)L^2+E_r(z)\Li_2(z).
\end{equation}
Thus the adjacent-two-pole dilogarithm phenomenon extends to arbitrarily
separated integer poles, a repeated double pole, and any polynomial part.
\end{corollary}
\begin{proof}
The possible nonempty pole profiles are $(1)$, $(1,1)$, and $(2)$.
Their stabilized weights are at most two. The positive-index words of
weight at most two are $\Li_1=L$, $\Li_2$, and
$\Li_{1,1}=L^2/2$. Apply Theorem~\ref{thm:compiler}.
\end{proof}

\begin{corollary}[Euler primitives and a finite differential module]
\label{thm:dmodule}
For every integer $e\ge0$,
\begin{equation}\label{eq:eulerprimitive}
 \sum_{M\ge1}\frac{D_{r,M}(R)}{M^e}z^M
             =\mathcal J_0^eG_r(z;R)\in\mathscr L_{W_r+e}.
\end{equation}
All generators $G_r$ for a fixed weight $R$ lie in a common differential
module over $K(z)$ spanned by at most $2^{W_\infty}$ functions.
Consequently each individual $G_r$ satisfies a nonzero scalar linear
differential equation over $K(z)$ of order at most $2^{W_\infty}$.
\end{corollary}
\begin{proof}
The first statement follows by anchored iteration. There are
$2^{w-1}$ compositions of positive weight $w$, so the empty word and
all words of weight at most $W_\infty$ give $2^{W_\infty}$ generators.
Differentiation maps their span into itself by the word derivative
rule and the product rule for rational factors. The first
$2^{W_\infty}+1$ derivatives of an individual element are therefore
linearly dependent over $K(z)$. This proves the scalar equation claim.
It does \emph{not} assert that the same scalar equation annihilates
all ranks; rational prefactors vary with $r$.
\end{proof}

\section{Several simple poles: fixed-size highest coefficients}
\label{sec:hahn}

For this section omit the polynomial part and write
\begin{equation}\label{eq:simpleR}
 R(x)=\sum_{j=1}^d\frac{c_j}{x+a_j}=\frac{P_*(x)}{L(x)},
 \quad L(x)=\prod_{j=1}^d(x+a_j),\quad
 P_*(x)=\sum_{j=1}^d c_j\frac{L(x)}{x+a_j},
\end{equation}
where $a_j$ are distinct nonnegative integers and $c_j\ne0$.
After \eqref{eq:shift}, let $X$ stand for $H_M$. The formal determinant
is a polynomial in $X$ with coefficients rational in $M$.

\begin{theorem}[The top coefficient is a polynomial-weight determinant]
\label{thm:compression}
For $r\ge d$, put $q=r-d$ and
$T_{r,M}=[X^d]D_{r,M}(R)$. Then
\begin{equation}\label{eq:top}
 \boxed{\displaystyle
 T_{r,M}=\left(\prod_{j=1}^d c_j\right)
          \prod_{i<j}(a_j-a_i)^2\,
     \det\left(\sum_{n=1}^M n^{i+j}L(n)P_*(n)\right)_{0\le i,j<q}.}
\end{equation}
The determinant of size zero is one. In particular the highest
harmonic coefficient has no rational cutoff poles.
\end{theorem}
\begin{proof}
The coefficient of $X$ in the moment matrix is
$\sum_j c_jv_jv_j^{\mathsf T}$, where
$v_j=(1,-a_j,\ldots,(-a_j)^{r-1})^{\mathsf T}$.
Use the monic polynomial basis
\[
       1,x,\ldots,x^{d-1},L(x),xL(x),\ldots,x^{q-1}L(x).
\]
Its triangular coefficient matrix has determinant one. Each of the
last $q$ basis vectors vanishes at every $-a_j$, so the harmonic
matrix has support only in the first $d$ rows and columns. Its
$d\times d$ determinant is
$\prod_j c_j\prod_{i<j}(a_j-a_i)^2$.
The coefficient of $X^d$ must select this entire block. The complementary
moment matrix sums $x^{i+j}L(x)^2R(x)=x^{i+j}L(x)P_*(x)$, a polynomial
weight. This proves \eqref{eq:top}, including its normalization.
\end{proof}

For completeness, the monic uniform Hahn polynomial on $1,\ldots,M$
can be written without hypergeometric denominators as
\begin{equation}\label{eq:hahn}
 p_j(x;M)=\binom{2j}{j}^{-1}
 \sum_{\ell=0}^j(-1)^{j+\ell}\binom j\ell\binom{j+\ell}{\ell}
       \fall{(M-1-\ell)}{j-\ell}\fall{(x-1)}\ell.
\end{equation}
For $j<M$ it is orthogonal with unit weights and has squared norm
\begin{equation}\label{eq:hahnnorm}
 h_j(M)=\frac{\Gamma(M+j+1)}
                  {(2j+1)\Gamma(M-j)\binom{2j}{j}^{\!2}}.
\end{equation}
These are the uniform specialization of the classical finite Hahn
formulas \cite{dlmf-hahn}. They imply
$\det(\sum_{n=1}^M n^{i+j})_{i,j<q}=\prod_{j<q}h_j(M)$ by a monic
triangular change of basis. The endpoint identity used in
Theorem~\ref{thm:sharp} follows by putting $x=0$ in \eqref{eq:hahn}
and applying the finite binomial Vandermonde identity.

\begin{corollary}[At most $2d-1$ Hahn rows, at every particle count]
\label{thm:fixedsize}
Let $W(x)=L(x)P_*(x)=\lambda\prod_{i=1}^e(x-b_i)$ have its actual
nonzero leading coefficient $\lambda$ and degree $e\le2d-1$.
For distinct roots and $M\ge q+e$,
\begin{equation}\label{eq:fixedsize}
 \begin{split}
 T_{r,M}={}&\left(\prod_jc_j\right)\prod_{i<j}(a_j-a_i)^2
      \lambda^q(-1)^{eq}\mathcal U_q(M)\\
 &\hspace{12mm}\times
       \frac{\det(p_{q+j-1}(b_i;M))_{1\le i,j\le e}}
                    {\prod_{i<j}(b_j-b_i)}.
 \end{split}
\end{equation}
Repeated roots are handled by the confluent divided-difference limit.
The compressed formula \eqref{eq:top} remains the definition at every
smaller cutoff and whenever the leading coefficient changes degree.
Thus the highest harmonic coefficient has a determinant of fixed size
at most $2d-1$, independent of $r$.
\end{corollary}
\begin{proof}
For a measure $\nu$ and its monic orthogonal polynomials, the
Christoffel determinant identity is
\[
 D_q\!\left[\prod_{i=1}^e(x-b_i)\nu\right]
 =(-1)^{eq}D_q[\nu]\,
           \frac{\det(p_{q+j-1}(b_i))_{i,j=1}^e}{\prod_{i<j}(b_j-b_i)}.
\]
One proof appends the $e$ fixed nodes $b_i$ to the Vandermonde in the
$q$ integration nodes. The cross factors are
$\prod_{i,j}(b_i-x_j)$, giving the stated sign. Replace monomials by
monic orthogonal polynomials and integrate the first $q$ rows against
the companion Vandermonde. Orthogonality leaves the product of the
first $q$ squared norms and the displayed $e$ rows. All steps are
finite for the discrete measure here. This is the standard
Christoffel transformation \cite{krattenthaler}. Applying it to
\eqref{eq:top}, and extracting the constant factor $\lambda^q$,
proves the result. Coalescing roots replaces evaluation rows by
derivative rows divided by the corresponding factorials, because the
Vandermonde quotient is a symmetric polynomial in the roots.
\end{proof}

\subsection{Exceptional residues as an exact matrix spectrum}

Let $M\ge q$ and $q>0$. For real $a_j\ge0$ define positive definite
$q\times q$ matrices
\begin{equation}\label{eq:Bres}
 B_j=\left(\sum_{n=1}^M
             n^{u+v}\frac{L(n)^2}{n+a_j}\right)_{0\le u,v<q}.
\end{equation}
For a nonzero polynomial $f$ of degree below $q$, its quadratic form
is a positive sum of $f(n)^2$ over at least $q$ distinct nodes. Hence
$B_j$ is positive definite.

\begin{theorem}[Residue hypersurface and real one-parameter slices]
\label{thm:stability}
Apart from the explicit factors $\prod_jc_j$ in \eqref{eq:top}, the
highest-coefficient exceptional locus is exactly
\begin{equation}\label{eq:exception}
           \det\left(\sum_{j=1}^d c_jB_j\right)=0.
\end{equation}
The polynomial on the left is nonzero when every $\operatorname{Im}c_j>0$,
and also when every $c_j>0$ is real. Fix positive real values for all
residues except $c_i=\tau$, and set
$A_i=\sum_{j\ne i}c_jB_j$. If $d\ge2$, then all exceptional values
of $\tau$ in \eqref{eq:exception} are
\begin{equation}\label{eq:roots}
                  \tau=-\lambda_1,\ldots,-\lambda_q,
\end{equation}
where $\lambda_j>0$ are the eigenvalues of
$B_i^{-1/2}A_iB_i^{-1/2}$, counted with multiplicity. No simplicity
assertion is made without additional hypotheses.
\end{theorem}
\begin{proof}
The compressed determinant in \eqref{eq:top} is exactly
$\det\sum_jc_jB_j$. If the imaginary parts are positive and a nonzero
complex vector $v$ were in its kernel, taking the imaginary part of
$v^*(\sum_jc_jB_j)v$ would give
$\sum_j\operatorname{Im}c_j\,v^*B_jv=0$, a contradiction. The positive
real case is immediate. This is the standard positive-matrix proof
of determinantal stability \cite{borcea}, specialized to the present
compressed measure. Finally
\[
 \det(A_i+\tau B_i)=\det B_i\,
                  \prod_{j=1}^q(\tau+\lambda_j),
\]
and $A_i$ is positive definite for $d\ge2$. This proves the exact
spectral description. The extra zero $\tau=0$ of the full top
coefficient changes the active pole set and is kept separate.
\end{proof}

\section{Exact endpoint cancellation equations}
\label{sec:endpoints}

The highest coefficient in \eqref{eq:top} is polynomial in $M$, but
lower harmonic coefficients can have cutoff poles. They must be
retained before forming a generating function. For the simple-pole
weight \eqref{eq:simpleR}, polynomial division gives the entirely
explicit moments
\begin{equation}\label{eq:simplemoment}
 \begin{split}
 \mu_j(M,X)=\sum_{a\in A}c_a\bigg\{&
  \sum_{v=0}^{j-1}(-a)^{j-1-v}S_v(M)\\
 &+(-a)^j\left(X+\sum_{b=1}^a\frac1{M+b}-H_a\right)\bigg\}.
 \end{split}
\end{equation}
The inner polynomial sum is empty if $j=0$.

For $1\le h\le\max A$, put
\[
 A_h^+=\{a\in A:a\ge h\},\qquad d_h=|A_h^+|,
 \qquad v_a=(1,-a,\ldots,(-a)^{r-1})^{\mathsf T}.
\]
Let $\boldsymbol\mu$ denote the formal moment matrix from
\eqref{eq:simplemoment}, and define its regular part at $M=-h$ by
\begin{equation}\label{eq:regularmatrix}
 B_h(M,X)=\boldsymbol\mu(M,X)-\frac1{M+h}
                         \sum_{a\in A_h^+}c_av_av_a^{\mathsf T}.
\end{equation}
Every entry is analytic in $M$ near $-h$ after its removable
cancellations are simplified.

\begin{theorem}[Every principal coefficient from a finite bordered jet]
\label{thm:endpoints}
The order of the formal cutoff pole at $M=-h$ is at most
$\min(r,d_h)$. Write its principal part as
$\sum_{\ell=1}^{\min(r,d_h)}P_{h,\ell}(X)/(M+h)^\ell$.
For a subset $S\subset A_h^+$ of size $k$, let $V_S$ have columns
$v_a$ for $a\in S$. Then
\begin{equation}\label{eq:endpointjet}
 \boxed{\displaystyle
 P_{h,\ell}(X)=
 \sum_{k=\ell}^{\min(r,d_h)}\ \sum_{\substack{S\subset A_h^+\\|S|=k}}
 \frac{(-1)^k\prod_{a\in S}c_a}{(k-\ell)!}
 \left.
 \partial_M^{k-\ell}
 \det\begin{pmatrix}B_h(M,X)&V_S\\V_S^{\mathsf T}&0\end{pmatrix}
 \right|_{M=-h}.}
\end{equation}
Thus the pole cancels identically as a formal harmonic expression if
and only if all of these polynomials in $X$ vanish. More generally,
its exact order is the largest $\ell$ with $P_{h,\ell}\ne0$.
Moreover
\begin{equation}\label{eq:endpointdegree}
              \deg_X P_{h,\ell}+\ell\le\min(r,d).
\end{equation}
\end{theorem}
\begin{proof}
For any square matrix $B$ and vectors $v_a$, the rank-one update
expansion is
\begin{equation}\label{eq:border}
 \det\left(B+\sum_a t_av_av_a^{\mathsf T}\right)
 =\sum_S(-1)^{|S|}\left(\prod_{a\in S}t_a\right)
          \det\begin{pmatrix}B&V_S\\V_S^{\mathsf T}&0\end{pmatrix}.
\end{equation}
For the empty subset the bordered determinant is $\det B$.
One can prove the identity first when $B$ is invertible by the
Schur complement and the principal-minor expansion of
$\det(I+\operatorname{diag}(t)V^{\mathsf T}B^{-1}V)$.
Both sides are polynomials in the entries of $B$, so the identity holds
without invertibility. Terms with $|S|>r$ vanish.

Apply \eqref{eq:border} with $B=B_h$ and $t_a=c_a/(M+h)$.
Taylor expansion of each analytic bordered determinant gives
\eqref{eq:endpointjet}. It also proves the maximal pole order and the
necessary-and-sufficient cancellation criterion. Finally, in the
simple-pole case the determinant is multiaffine in the original
$h_{a,1}$ and contains at most $\min(r,d)$ such factors. Selecting
$\ell$ inverse-cutoff factors leaves at most $\min(r,d)-\ell$ harmonic
factors. Equivalently use Lemma~\ref{thm:cutoffnormal} with all
$s_a=1$. This proves \eqref{eq:endpointdegree}.
\end{proof}

This is an exact finite classification by explicit polynomial equations,
not an intrinsic formula for the first particle count at which every
endpoint disappears. Such a threshold, expressed only in pole locations
and residues without constructing the jets, remains a further question.
The theorem nevertheless gives every endpoint contribution required by
the compiler, including cases in which a naive residue-rank estimate is
strict.

\subsection{A complete three-pole example}

Take
\begin{equation}\label{eq:threeR}
          R(n)=\frac1n+\frac1{n+1}+\frac1{n+2},\qquad r=3.
\end{equation}
The exact formal determinant has the following compact
polynomial-plus-principal-part representation:
\begin{equation}\label{eq:threeformal}
 \begin{split}
 D_{3,M}(R)={}&4H_M^3+A_2(M)H_M^2+A_1(M)H_M+A_0(M)\\
 &+\frac{8H_M^2-36H_M-19/2}{M+1}
   +\frac{4H_M-18}{(M+1)^2}
   +\frac{4H_M^2-26H_M+13}{M+2},
 \end{split}
\end{equation}
where
\begin{align}
 A_2(M)={}&\tfrac92M^4+27M^3+\tfrac{111}2M^2+45M-10,\notag\\
 A_1(M)={}&\tfrac38M^6+\tfrac{27}8M^5-\tfrac{45}4M^4
       -\tfrac{753}8M^3-\tfrac{1199}8M^2-\tfrac{69}2M+45,\notag\\
 A_0(M)={}&-\tfrac{23}{16}M^6-\tfrac{129}{16}M^5
       +\tfrac{219}{16}M^4+\tfrac{1527}{16}M^3
       +\tfrac{495}8M^2-\tfrac{169}4M+21.
 \label{eq:threepolys}
\end{align}
These formulas follow by inserting \eqref{eq:threeR} in
\eqref{eq:simplemoment} and expanding a $3\times3$ determinant;
the accompanying exact checks also compare them to the original
finite sums. The top coefficient is $4$, precisely
$\prod_{i<j}(a_j-a_i)^2$ at $a_j=0,1,2$. The double cutoff pole at
$M=-1$ shows why retaining only simple endpoint poles is incorrect for
three simple summand poles. This is a counterexample to that possible
extrapolation, not a claim that the adjacent-pole source made it.

Let $F_j(z)=\sum_{M\ge1}H_M^jz^M$, with $F_0=z/(1-z)$.
The finite product rules give
\begin{align}
 F_1&=\frac{\Li_1(z)}{1-z},\notag\\
 F_2&=\frac{\Li_2(z)+2\Li_{1,1}(z)}{1-z},\notag\\
 F_3&=\frac{\Li_3(z)+3\Li_{2,1}(z)+3\Li_{1,2}(z)
                           +6\Li_{1,1,1}(z)}{1-z}.
 \label{eq:F123}
\end{align}
An explicit identity for the three-pole generator is therefore
\begin{equation}\label{eq:threegen}
 \begin{split}
 G_3(z;R)={}&4F_3+A_2(\theta)F_2+A_1(\theta)F_1+A_0(\theta)F_0\\
 &+\mathcal J_1(8F_2-36F_1-\tfrac{19}2F_0)
    +\mathcal J_1^2(4F_1-18F_0)\\
 &+\mathcal J_2(4F_2-26F_1+13F_0).
 \end{split}
\end{equation}
Here each $\mathcal J_h$ is the explicitly anchored integral
\eqref{eq:inverse}. The supplied compiler removes them completely and
returns eight rationally weighted words of length at most three. The
fully expanded coefficients are recorded in
\src{data/compiled_examples.json}; \eqref{eq:threegen} is the more
readable exact expression. The derivative operators are of finite
order and act on the entire functions immediately following them.

\section{A small exact identity with a rank-independent extension}

\subsection{The triple pole and its weight-four primitive}

For the triple pole at the origin and two particles, put
\begin{equation}\label{eq:KM}
 K_M=H_M^{(1)}H_M^{(3)}-\bigl(H_M^{(2)}\bigr)^2.
\end{equation}
The diagonal terms in the two quasi-shuffles cancel. Consequently
\begin{equation}\label{eq:Kgen}
 \boxed{\displaystyle
 \sum_{M\ge1}K_Mz^M=
 \frac{\Li_{3,1}(z)+\Li_{1,3}(z)-2\Li_{2,2}(z)}{1-z}.}
\end{equation}
This concrete weight-four identity is also the first nontrivial member
of the spectral family in the next section. An anchored primitive is
\begin{equation}\label{eq:Kprimitive}
 \int_0^z\frac{1-t}{t}\sum_{M\ge1}K_Mt^M\dd t
           =\Li_{4,1}(z)+\Li_{2,3}(z)-2\Li_{3,2}(z).
\end{equation}
The right side is single-valued on the disk under the power-series
branch convention. There is no additive constant because both sides
vanish at zero.

\begin{proposition}[An ordinary convergent zeta evaluation]
\label{thm:Kendpoint}
\begin{equation}\label{eq:Kendpoint}
 \boxed{\displaystyle
 \sum_{M=1}^{\infty}
 \frac{H_M^{(1)}H_M^{(3)}-(H_M^{(2)})^2}{M(M+1)}
           =\frac{35}{2}\zeta(5)-9\zeta(2)\zeta(3).}
\end{equation}
\end{proposition}
\begin{proof}
The numerator is $O(1+\log M)$, so the sum is absolutely convergent.
Alternatively \eqref{eq:ensemble} expresses it as
$\sum_{n>m\ge1}(n-m)^2/(n^3m^3)$ up to the cutoff; all terms are
nonnegative. At $z=1$, integrating \eqref{eq:Kgen} as in
\eqref{eq:Kprimitive} therefore gives
\begin{equation}\label{eq:Kzetas}
                  \zeta(4,1)+\zeta(2,3)-2\zeta(3,2).
\end{equation}
Here $\zeta(a,b)=\sum_{n>m\ge1}n^{-a}m^{-b}$.

For transparency, the classical weight-five reduction can be checked
with three identities. Write $P=\zeta(2)\zeta(3)$ and $Z=\zeta(5)$.
The height-one Gamma identity gives
$\zeta(4,1)=2Z-P$. One way to see its exact coefficient is to use
\begin{equation}\label{eq:heightone}
 1-\frac{\Gamma(1-x)\Gamma(1-y)}{\Gamma(1-x-y)}
 =\sum_{m,n\ge1}\zeta(m+1,\{1\}^{n-1})x^my^n.
\end{equation}
For a direct normalization check, sum the height-one iterated-integral
formula over $m,n$ to obtain
\[
 y\int_0^1(t^{-x}-1)(1-t)^{-y-1}\dd t
   =y\{B(1-x,-y)-B(1,-y)\}
   =1-\frac{\Gamma(1-x)\Gamma(1-y)}{\Gamma(1-x-y)}.
\]
The difference integral is analytic for small $x,y$; the separated beta
terms may first be evaluated with $\operatorname{Re}y<0$ and then
continued. This is consistent with the height-one Gamma convention in
the repository \cite{harmonic}. Its logarithmic
Gamma expansion gives coefficient $x^3y^2$ equal to
$2\zeta(5)-\zeta(2)\zeta(3)$: the linear exponent contributes $2\zeta(5)$
and the product of its weight-two and weight-three terms contributes
$-\zeta(2)\zeta(3)$.

The finite-sum stuffle, passed to the absolutely convergent limit, gives
$P=\zeta(2,3)+\zeta(3,2)+Z$. The iterated-integral shuffle of words
$01$ and $001$ gives
$P=\zeta(2,3)+3\zeta(3,2)+6\zeta(4,1)$; the multiplicities follow by
listing the ten order-preserving interleavings. Solving yields
\[
 \zeta(3,2)=3P-\tfrac{11}{2}Z,\qquad
 \zeta(2,3)=\tfrac92Z-2P.
\]
Substitute in \eqref{eq:Kzetas}. This proves \eqref{eq:Kendpoint} by
exact identities, without a numerical-relation search.
\end{proof}

The reduction of \eqref{eq:Kzetas} is classical Euler-sum algebra. The
structural extension here is that every determinant
$D_{r,M}(n^{-3})$, for arbitrary $r$, admits a generator in the same
weight-four module, with the explicit nonzero coefficient
\eqref{eq:sharp}. For example the next determinant is
\begin{equation}\label{eq:triple-r3}
 \begin{split}
 D_{3,M}(n^{-3})={}&S_1(M)
       \{H_M^{(1)}H_M^{(3)}-(H_M^{(2)})^2\}\\
 &-M^2H_M^{(3)}+2M H_M^{(1)}H_M^{(2)}-(H_M^{(1)})^3.
 \end{split}
\end{equation}
Its generator follows by quasi-shuffling just the displayed products
and replacing $M$ by $\theta$. Higher rank changes the rational
coefficients and differential operators, not the maximal word length.

\subsection{A shifted double pole with an ordinary zeta endpoint}

The repeated-pole case of Corollary~\ref{thm:quadratic} gives another
small identity. Set
\[
 B_M=M\{H_{M+1}^{(2)}-1\}-(H_{M+1}-1)^2,
 \qquad L=-\log(1-z).
\]
A monic change from the basis $1,n$ to $1,n+1$ shows that $B_M$ is
$D_{2,M}((n+1)^{-2})$.

\begin{proposition}[Shifted double-pole generator and Euler primitive]
\label{thm:shifteddouble}
For $|z|<1$,
\begin{equation}\label{eq:doublegen}
 \boxed{\displaystyle
 \sum_{M\ge1}B_Mz^M=
 \frac{(3z-2)\Li_2(z)+(1-z)(3L-L^2)-z}{z(1-z)^2}.}
\end{equation}
The value at zero is removable. Moreover
\begin{equation}\label{eq:doubleendpoint}
 \boxed{\displaystyle
 \sum_{M\ge1}\frac{M\{H_{M+1}^{(2)}-1\}-(H_{M+1}-1)^2}
                       {M(M+1)(M+2)}
       =\frac52\zeta(3)-\frac32\zeta(2)-\frac12.}
\end{equation}
\end{proposition}
\begin{proof}
Let $U=\sum_{M\ge1}(H_{M+1}^{(2)}-1)z^M$ and
$V=\sum_{M\ge1}(H_{M+1}-1)^2z^M$. The elementary shift of
\eqref{eq:hgen} and the quadratic quasi-shuffle give
\[
 U=\frac{\Li_2(z)}{z(1-z)}-\frac1{1-z},\qquad
 V=\frac{\Li_2(z)+L^2-2L}{z(1-z)}+\frac1{1-z}.
\]
Their combination $\theta U-V$ is precisely \eqref{eq:doublegen}.

The determinant is a nonnegative finite Vandermonde sum and is $O(M)$.
Its endpoint series is therefore absolutely convergent. Since
\[
 \frac12\int_0^1z^{M-1}(1-z)^2\dd z=\frac1{M(M+1)(M+2)},
\]
it equals $\tfrac12\int_0^1N(z)/z^2\dd z$, where
$N(z)=(3z-2)\Li_2(z)+(1-z)(3L-L^2)-z$.
Here $N(z)/z\to0$ at zero and $N(1)=\zeta(2)-1$. Direct
differentiation gives
\[
 N'(z)=3\Li_2(z)+L^2-2L+2-\frac{2L}{z}.
\]
Integrate by parts once. Use
$\int_0^1\Li_2(z)\dd z/z=\zeta(3)$,
$\int_0^1L\dd z/z=\zeta(2)$, and
$\int_0^1L^2\dd z/z=2\zeta(2,1)=2\zeta(3)$.
The last equality follows also from coefficient $xy^2$ in
\eqref{eq:heightone}. The remaining terms must be kept together:
\[
 \int_0^1\left(\frac1z-\frac{L}{z^2}\right)\dd z
       =\left[\frac Lz-L\right]_0^1=-1.
\]
Thus the integral is
$\tfrac12\{-\zeta(2)+1+5\zeta(3)-2\zeta(2)-2\}$,
which proves \eqref{eq:doubleendpoint}.
\end{proof}

\section{A spectral master identity at every rank}
\label{sec:spectral}

The integer-pole theorem has a useful analytic companion that should
not be confused with differentiating its formal weight bound. Define,
for arbitrary complex $s$,
\begin{equation}\label{eq:spectraldet}
 \mathcal K_{r,M}(s)=
            \det\bigl(H_M^{(s-i-j)}\bigr)_{0\le i,j<r}.
\end{equation}
For a finite cutoff, each entry is an entire function of $s$, with
$n^{-s}=\exp(-s\log n)$ using the real logarithm. In this section
$\Li_{a_1,\ldots,a_r}(z)$ is also allowed to have arbitrary complex
indices; its defining nested series is still entire in all indices
when $|z|<1$.

Let $\mathfrak S_r$ act on $\{0,\ldots,r-1\}$. For
$\sigma,\tau\in\mathfrak S_r$, put
\begin{equation}\label{eq:asig}
             a_j(s;\sigma,\tau)=s-\sigma(j)-\tau(j),
                     \qquad 0\le j<r.
\end{equation}

\begin{theorem}[All-rank spectral polylogarithm and zeta identities]
\label{thm:spectral}
For every $s\in\C$ and $|z|<1$,
\begin{equation}\label{eq:spectralgen}
 \boxed{\displaystyle
 \sum_{M\ge1}\mathcal K_{r,M}(s)z^M
 =\frac1{1-z}\sum_{\sigma,\tau\in\mathfrak S_r}
       \operatorname{sgn}(\sigma)\operatorname{sgn}(\tau)
       \Li_{a_0(s),\ldots,a_{r-1}(s)}(z).}
\end{equation}
Both sides are entire in $s$ and can be differentiated to every order
on compact subsets. If $\operatorname{Re}s>2r-2$, then the ordinary
endpoint sum is absolutely convergent and equals
\begin{equation}\label{eq:spectralendpoint}
 \boxed{\displaystyle
 \sum_{M\ge1}\frac{\mathcal K_{r,M}(s)}{M(M+1)}
 =\sum_{\sigma,\tau\in\mathfrak S_r}
       \operatorname{sgn}(\sigma)\operatorname{sgn}(\tau)
       \zeta(a_0(s)+1,a_1(s),\ldots,a_{r-1}(s)).}
\end{equation}
The endpoint identity is holomorphic in this half-plane and also
admits all fixed-order spectral derivatives there.
\end{theorem}
\begin{proof}
Apply \eqref{eq:ensemble} with $R(n)=n^{-s}$. Reverse the node order
to $n_0>\cdots>n_{r-1}$, which does not change the square of the
Vandermonde. Expand that square as two determinants:
\[
 \Delta(n_0,\ldots,n_{r-1})^2
 =\sum_{\sigma,\tau}\operatorname{sgn}(\sigma)
       \operatorname{sgn}(\tau)\prod_{j=0}^{r-1}
                           n_j^{\sigma(j)+\tau(j)}.
\]
Sum the cutoff $M\ge n_0$ geometrically. This gives
\eqref{eq:spectralgen}, including the absence of an $r!$ factor.
For $s$ on a compact set, all finite harmonic moments and their fixed
spectral derivatives have polynomial growth in $M$ times powers of
$\log(2+M)$. The factor $z^M$ supplies a locally uniform summable
majorant, proving entireness and every fixed-order differentiation.
The same argument applies to each nested polylogarithm.

For the endpoint, the cutoff tail is exactly
\[
             \sum_{M=n_0}^{\infty}\frac1{M(M+1)}=\frac1{n_0}.
\]
To justify all absolute rearrangements in the claimed half-plane,
use the elementary rearrangement inequality on the decreasing nodes:
for every permutation $\sigma$,
$\prod_jn_j^{\sigma(j)}\le\prod_jn_j^{r-1-j}$.
It follows by successively swapping any inverted pair of exponents.
After both Vandermonde expansions and the extra factor $1/n_0$, every
absolute summand is dominated by
\[
 n_0^{-\operatorname{Re}s+2r-3}
             \prod_{j=1}^{r-1}n_j^{-\operatorname{Re}s+2r-2-2j}.
\]
When $\operatorname{Re}s>2r-2$, all the positive decay exponents in
this expression exceed one. Dropping the ordering gives a product of
convergent ordinary zeta series. A compact sub-half-plane permits
additional fixed powers of logarithms by slightly decreasing those
exponents. This proves convergence, \eqref{eq:spectralendpoint}, and
all stated derivatives.
\end{proof}

At a positive integer $s$, \eqref{eq:spectralgen} can contain
nonpositive indices if $r$ is large. Theorem~\ref{thm:compiler} removes
them constructively and yields the uniform alphabet of maximal weight
$\lfloor(s+1)^2/4\rfloor$, independent of $r$. The unreduced spectral
formula is not itself a minimal-weight representation. At a general
noninteger $s$, or after order differentiation, the finite integer-weight
alphabet theorem is not being asserted.

\section{Generalized Stieltjes constants and exact spectral jets}
\label{sec:stieltjes}

Use the Laurent convention
\begin{equation}\label{eq:stieltjes}
 \zeta(1+u,a)=\frac1u+
          \sum_{q\ge0}\frac{(-1)^q}{q!}\gamma_q(a)u^q.
\end{equation}
For $p\in\C$, define the finite logarithmic harmonic sums
\begin{equation}\label{eq:Lpq}
 L_{q,p}(M)=\sum_{n=1}^M\frac{\log^q n}{n^p}
                  =(-1)^q\partial_p^qH_M^{(p)}.
\end{equation}
The Hurwitz difference $H_M^{(p)}=\zeta(p)-\zeta(p,M+1)$ is entire
in $p$ because its two residues at $p=1$ cancel. Comparing
\eqref{eq:stieltjes} gives the ordinary finite identity
\begin{equation}\label{eq:Jq}
          J_q(M):=\gamma_q(1)-\gamma_q(M+1)=L_{q,1}(M).
\end{equation}
This is a difference of Stieltjes constants, not a value assigned to
one of the separately singular terms $\zeta^{(q)}(1,a)$.
For $p\ge2$ integral,
\begin{equation}\label{eq:Lzeta}
 L_{q,p}(M)=(-1)^q
                  \{\zeta^{(q)}(p)-\zeta^{(q)}(p,M+1)\}.
\end{equation}
At $q=0$ the same harmonic variables are polygamma differences:
\begin{equation}\label{eq:polyGamma}
 H_M^{(p)}=\frac{(-1)^{p-1}}{(p-1)!}
       \{\psi^{(p-1)}(M+1)-\psi^{(p-1)}(1)\},
                   \qquad p\ge1.
\end{equation}
For $p=1$ this uses $\psi^{(0)}=\psi$ and is the usual digamma shift.
These standard Hurwitz and polygamma identities are consistent with
DLMF conventions \cite{dlmf-zeta}; the finite sums above also prove
the needed cancellations directly.

\subsection{A complete rank-two Stieltjes family}

Define a holomorphic function for $\operatorname{Re}s>2$ by
\begin{equation}\label{eq:Zmaster}
 \mathcal Z(s)=\zeta(s+1,s-2)+\zeta(s-1,s)-2\zeta(s,s-1).
\end{equation}
The double-zeta series in this formula all converge absolutely on
that half-plane.

\begin{theorem}[All-order Stieltjes--harmonic endpoint identity]
\label{thm:stieltjes}
For every integer $q\ge0$,
\begin{equation}\label{eq:allstieltjes}
 \boxed{\displaystyle
 \sum_{M=1}^{\infty}\frac1{M(M+1)}
    \sum_{j=0}^q\binom qj
       \bigl\{L_{j,3}(M)J_{q-j}(M)
                   -L_{j,2}(M)L_{q-j,2}(M)\bigr\}
                   =(-1)^q\mathcal Z^{(q)}(3).}
\end{equation}
Every series here is an ordinary absolutely convergent series. At
$q=0$ it is \eqref{eq:Kendpoint}. For $|z|<1$ the corresponding
cutoff generator is
\begin{equation}\label{eq:stieltjesgen}
 \begin{split}
 &\sum_{M\ge1}z^M\sum_{j=0}^q\binom qj
       \{L_{j,3}(M)J_{q-j}(M)-L_{j,2}(M)L_{q-j,2}(M)\}\\
 &\qquad=
 \frac{(-1)^q}{1-z}
 \left.\partial_s^q
  \{\Li_{s,s-2}(z)+\Li_{s-2,s}(z)-2\Li_{s-1,s-1}(z)\}
                                                      \right|_{s=3}.
 \end{split}
\end{equation}
\end{theorem}
\begin{proof}
For $r=2$, \eqref{eq:spectraldet} is
$H_M^{(s)}H_M^{(s-2)}-(H_M^{(s-1)})^2$.
Differentiate this finite expression $q$ times at $s=3$, multiply by
$(-1)^q$, and use Leibniz's rule, \eqref{eq:Lpq}, and \eqref{eq:Jq}.
The result is precisely the inner binomial sum. Apply the corresponding
derivative to \eqref{eq:spectralendpoint} and
\eqref{eq:spectralgen}. Theorem~\ref{thm:spectral} justifies both
interchanges. Directly, $L_{j,2}$ and $L_{j,3}$ remain bounded in $M$
for each fixed $j$, whereas $J_j(M)=O(1+\log^{j+1}(2+M))$. This also
checks absolute convergence of the displayed endpoint series.
\end{proof}

For example, the first logarithmic identity is
\begin{equation}\label{eq:qone}
 \sum_{M\ge1}\frac{
    L_{1,3}(M)H_M^{(1)}+H_M^{(3)}\{\gamma_1(1)-\gamma_1(M+1)\}
          -2H_M^{(2)}L_{1,2}(M)}{M(M+1)}=-\mathcal Z'(3).
\end{equation}
This is an exact relation between generalized Stieltjes constants,
Hurwitz-zeta derivatives, and depth-two zeta derivatives. It is not
obtained by differentiating the isolated evaluation
$\mathcal Z(3)=\tfrac{35}2\zeta(5)-9\zeta(2)\zeta(3)$.
That evaluation holds at one value of $s$; it does not determine
$\mathcal Z'(3)$ or higher derivatives. Promoting a special-value
identity to a parameter identity without proof would give an invalid
reduction of \eqref{eq:qone}.

\subsection{The all-rank version at the Stieltjes boundary}

Let $s_0=2r-1$ and $q\ge0$. All integer indices
$s_0-i-\sigma(i)$ in the following formula are positive. Expanding the
determinant and its derivative gives
\begin{equation}\label{eq:rankjet}
 \begin{split}
 \mathcal T_{r,q}(M)=
 \sum_{\sigma\in\mathfrak S_r}\operatorname{sgn}(\sigma)
 \sum_{q_0+\cdots+q_{r-1}=q}
       \frac{q!}{q_0!\cdots q_{r-1}!}
       \prod_{i=0}^{r-1}L_{q_i,s_0-i-\sigma(i)}(M).
 \end{split}
\end{equation}
With $a_j(s)=a_j(s;\sigma,\tau)$ as in \eqref{eq:asig},
Theorem~\ref{thm:spectral} proves
\begin{equation}\label{eq:rankstieltjes}
 \begin{split}
 \sum_{M\ge1}\frac{\mathcal T_{r,q}(M)}{M(M+1)}
 ={}&(-1)^q\left.\partial_s^q\left\{
 \sum_{\sigma,\tau\in\mathfrak S_r}
       \operatorname{sgn}(\sigma)\operatorname{sgn}(\tau)\right.\right.\\
 &\left.\left.\hspace{4mm}\times
 \zeta(a_0(s)+1,a_1(s),\ldots,a_{r-1}(s))
 \right\}\right|_{s=s_0}.
 \end{split}
\end{equation}
Formula \eqref{eq:Jq} replaces every index-one factor by an actual
Stieltjes difference. The other factors are \eqref{eq:Lzeta}; hence
this is a finite all-rank, all-order specification of ordinary
Stieltjes--zeta identities. There is at most one index-one factor in
any product: it requires $i=\sigma(i)=r-1$. At $q=0$, the zeta terms
on the right have total index weight $r^2+1$. Their possible further
reduction to a chosen multiple-zeta basis is a distinct arithmetic
problem, not needed for the identity.

\subsection{Anchored spatial primitives and the Gamma endpoint}

There is also a useful exact spatial primitive, stated here to fix its
normalization rather than as a new special-function fact. For $a>-1$,
let $J_q(M,a)=\gamma_q(a+1)-\gamma_q(M+a+1)$. For
$-1<\alpha<\beta$,
\begin{equation}\label{eq:spatialprimitive}
 \begin{split}
 \int_\alpha^\beta J_q(M,a)\dd a
 &=\frac1{q+1}\sum_{n=1}^M
            \{\log^{q+1}(n+\beta)-\log^{q+1}(n+\alpha)\}\\
 &=\frac{(-1)^{q+1}}{q+1}\left[
    \zeta^{(q+1)}(0,a+1)-\zeta^{(q+1)}(0,M+a+1)
                         \right]_{a=\alpha}^{a=\beta}.
 \end{split}
\end{equation}
The proof integrates each finite logarithmic summand and differentiates
the entire finite Hurwitz difference at zero. At $q=0$ it reduces to
\begin{equation}\label{eq:gammalogprimitive}
 \int_\alpha^\beta\{\psi(M+a+1)-\psi(a+1)\}\dd a
       =\left[\log\frac{\Gamma(M+a+1)}{\Gamma(a+1)}
                                                 \right]_\alpha^\beta.
\end{equation}
Positive real arguments fix the logarithm branch. These formulas may
also be continued to a simply connected complex domain avoiding all
$-1,-2,\ldots,-M$, with one consistent logarithm branch.

\section{Verification, claim status, and source corrections}
\label{sec:verification}

\subsection{Exact checks and numerical diagnostics have different roles}

The all-rank proofs above do not depend on fitting coefficient tables.
The included programs independently check finite instances and the
translation from formal identities to code. Their output is retained
in machine-readable JSON and ordinary run logs.

\begin{center}
\begin{tabular}{p{0.66\textwidth}r}
\toprule
Executed exact check & Count\\\midrule
Direct determinant versus compiled Taylor coefficient comparisons & 99\\
Single-pole weighted-degree checks, orders $1$--$7$, ranks $1$--$4$ & 28\\
Extremal Gamma-product coefficient checks on the same grid & 28\\
Several-simple-pole top-coefficient compression checks & 8\\
Local confluent matrix factorizations using independent polynomial division & 12\\
Endpoint bordered-jet principal coefficients & 7\\
All-rank double-permutation spectral normalizations & 24\\
Deliberately corrupted identities rejected & 2\\\bottomrule
\end{tabular}
\end{center}

The compiler checks include separated simple poles, three simple poles,
a shifted double pole, a mixture of double and simple poles, and a
polynomial part. For each of nine examples, the constant coefficient
and coefficients through $z^{10}$ are compared exactly with the
original finite-sum determinant. The maximum output word length is
checked against $W_r$. The local factorization tests use polynomial
long division of the rational summand rather than the compiler's
shifted-binomial formula. The endpoint tests compare the bordered-jet
formula to a separately expanded Laurent series.

The numerical suite uses 60 decimal digits and contains 28 diagnostics:
original cutoff determinants against the compiled expressions;
nonintegral and complex spectral parameters; Stieltjes differences
through order three; four differentiated spectral generators; and
independent Mellin integrals for the weight-five double-zeta reduction.
The suite uses a stated scaled tolerance of $10^{-48}$. Recorded residuals,
working precision, and truncation settings are in
\src{data/numeric_checks.json}. These are floating-point diagnostics,
not interval enclosures. A small residual alone is not used to prove
any theorem, and no proof-assistant formalization is claimed.

\subsection{What the package implements}

\src{verification/compiler.py} implements the rational moment formula,
finite quasi-shuffles, the two-letter integration recursion, and the
anchored inverse operators. Its output is a dictionary of binary
words and rational coefficients; an empty word means the rational
term. The example file records this complete representation, not just
numerical values. The finite algorithm works at arbitrary specified
rank, but no polynomial-time complexity claim is made: symbolic
expression size and partial-fraction expansion can grow substantially.

A minimal call is
\begin{verbatim}
from verification.compiler import compile_hankel, encode
result = compile_hankel(2, [(0, 3, 1)])
print(encode(result))
\end{verbatim}
Each triple is \texttt{(nonnegative integer shift, pole order, residue)}.
The output here is exactly \eqref{eq:Kgen} in binary-word notation.
The program is intended for exact rational or symbolic coefficients;
passing binary floating-point residues changes the nature of its
arithmetic and is not part of the exact certificate claims.

\subsection{Corrections and integration decisions}

No false mathematical theorem was found in the particular inherited
proof dependencies used here. The following are status corrections
and precision safeguards, not accusations of errors in those theorems.

The gap-six report's Question~Q4 still asks for the Gamma normalization
constants $K_m$. The newer incoming Gamma-normalization package gives
their explicit evaluation \cite{normalizations}; Q4 should point to
that result instead of remaining an unspecialized open question.
Similarly, the older mixed Gamma normal-form coefficient question is
already settled by the universal Gamma archive \cite{universal} and
should not be reopened on the strength of an older manuscript snapshot.

For the several-pole program, Questions~Q7--Q9 can now reference the
following specific advances: a fixed-size highest-coefficient formula
for any finite simple pole set; exact finite endpoint cancellation
jets; a uniform finite polylogarithm compiler including repeated
integer poles; and the spectral Euler primitives proved here.
A full intrinsic classification of all cutoff cancellations, or a
necessary-and-sufficient criterion for \emph{minimal classical}
polylogarithm depth, remains open. Those broader questions should not
be marked completely solved merely because a finite compiler exists.

Three conventions must survive integration. First, inverse-cutoff
operators are anchored at zero, and principal parts are retained before
summation. Second, all degree and weight statements concern formal
harmonic coordinates; they do not assert independence of evaluated
constants. Third, order derivatives of zeta functions are taken in
actual parameter identities, not in an isolated special-value
reduction. These distinctions are mathematically necessary.

\section{Further research questions}
\label{sec:questions}

\begin{enumerate}[label=\textbf{Q\arabic*.},leftmargin=*]
\item \textbf{Determine minimal functional weight and depth.}
Theorem~\ref{thm:sharp} proves formal sharpness, not transcendental
minimality. Which pole profiles force genuinely depth-two or higher
functions after all functional identities and argument changes are
allowed? In particular, classify the all-rank weights whose generators
use only classical $\Li_j$ and logarithms.

\item \textbf{Find a generic multiple-pole sharpness theorem.}
For a single pole the extremal coefficient is explicit. For several
poles, characterize when the optimized confluent minors survive their
compression and the merger of integer-shift harmonic variables.
A coefficient formula indexed by the maximizing allocations $k_a$
may distinguish generic sharpness from special residue cancellations.

\item \textbf{Replace endpoint jets by intrinsic cutoff invariants.}
Formula \eqref{eq:endpointjet} is a complete finite test, but it still
uses matrices depending on $r$. Seek exact vanishing thresholds and
factorizations directly from the integer dispersions of the poles,
their residues, and the polynomial part. The three-pole double
principal part in \eqref{eq:threeformal} is a first test case.

\item \textbf{Study the exceptional residue hypersurface.}
Equation \eqref{eq:exception} describes it exactly and gives its
positive-matrix slices. Determine its irreducible components,
singular loci, and conditions for simple roots on a one-residue slice.
Multiplicity cannot be discarded from the present theorem without a
further argument.

\item \textbf{Extend the compiler to rational and nonintegral shifts.}
For rational shifts, a roots-of-unity alphabet or a suitable cover is
natural. Derive a precise analogue of $W_r$ that tracks both the
confluent multiplicities and the required alphabet. Arbitrary complex
shifts should instead lead to a Lerch-type iterated-sum calculus.

\item \textbf{Reduce the Stieltjes spectral derivatives arithmetically.}
The identities \eqref{eq:allstieltjes} and \eqref{eq:rankstieltjes}
leave explicit derivatives of multiple zeta functions. Seek genuine
parameter identities reducing some of them to ordinary zeta or
Hurwitz-zeta derivatives. Differentiating the weight-five evaluation
at a single point is not an admissible shortcut.

\item \textbf{Make the differential module small and explicit.}
The finite $2^{W_\infty}$ generating set need not be minimal. Construct
cyclic equations, differential relations, or a reduced unipotent
system adapted to a fixed rational weight, and compare its rank
across residue specializations. The target is an exact common module,
not a falsely claimed single scalar equation for every particle count.

\item \textbf{Understand confluence at the level of identities.}
Continuous merging of simple poles differentiates shift parameters
and can introduce higher-weight harmonic variables. Explain precisely
how the simple-pole identities pass to the repeated-pole capacities,
including the passage from three simple poles to the weight-four
triple-pole sector, while keeping the endpoint subtractions uniform.
\end{enumerate}

\section*{Conclusion}
\addcontentsline{toc}{section}{Conclusion}
A finite pole set supplies more than a finite harmonic rank. Its
confluent derivative blocks determine an explicit weight allocation,
and that allocation survives every rational cutoff correction needed
for summation. This yields a finite, reproducible polylogarithm identity
at every rank. The single-pole coefficient \eqref{eq:sharp} proves the
formal invariant is attained and identifies its Gamma normalization.
The fixed-size compression, endpoint jets, and spectral Stieltjes
identities give complementary ways to extract exact consequences.
They extend a concrete current repository question without confusing
classical tools, newly proved combinations, finite evidence, and the
larger arithmetic problems that remain.

\begin{thebibliography}{99}
\small
\setlength{\itemsep}{3pt}

\bibitem{repo}
V. Reshetnikov, \emph{ProveIt}, polylogarithm manuscript and incoming
research packages. Repository observation at commit
\src{505c62fddea6b4a9eb61fee0c76e6c021e9c7128}, 11 October 2026 local time.
\url{https://github.com/VladimirReshetnikov/ProveIt}.

\bibitem{gap}
ProveIt incoming research report, \emph{Gap Six Harmonic Determinants},
archive \src{docs/incoming/ProveIt_Gap_Six_Harmonic_Determinants_2026-10-12.zip}.
In particular \src{sections/02_complex_rank.tex},
\src{sections/03_hahn_and_polylogs.tex}, and
\src{sections/05_audit_and_questions.tex}. Git blob
\src{d7695b94c3117d7bdcb5e75fcf930c410e9c31b6} at the comparison commit.

\bibitem{universal}
ProveIt incoming research report, \emph{A Universal Beta Formula for
Gamma Normal Forms}, archive
\src{docs/incoming/ProveIt_Universal_Gamma_Normal_Forms_2026-10-11.zip}.
README and result scope inspected. Git blob
\src{b2a5e51a1b340683721bb803c20410fc3124d9de}.

\bibitem{normalizations}
ProveIt incoming research report, \emph{Gamma Determinants and Colored
Stieltjes Moments}, archive
\src{docs/incoming/ProveIt_Gamma_Normalizations_and_Colored_Moments_2026-10-12.zip}.
Gamma determinant and normalization section inspected. Git blob
\src{a7d997b5d80ab90623ca58864871eb3829089628}.

\bibitem{harmonic}
ProveIt, \emph{Stieltjes harmonic identities}, harmonic Gamma-ratio and
height-one zeta development, source
\src{Analysis/Polylogarithms/docs/reports/stieltjes-harmonic-resonance/continuations/stieltjes-harmonic-identities/sections/harmonic.tex}.
Inspected Git blob \src{51108e8519d46af87dc3682956511fc525cecec3}.

\bibitem{hoffman}
M. E. Hoffman, \emph{Quasi-shuffle products}, Journal of Algebraic
Combinatorics \textbf{11} (2000), 49--68.
\url{https://arxiv.org/abs/math/9907173}.

\bibitem{krattenthaler}
C. Krattenthaler, \emph{Hankel determinants of linear combinations of
moments of orthogonal polynomials, II}, The Ramanujan Journal
\textbf{61} (2023), 597--627.
\url{https://doi.org/10.1007/s11139-021-00514-8}.

\bibitem{panzer}
E. Panzer, \emph{Algorithms for the symbolic integration of
hyperlogarithms with applications to Feynman integrals}, Computer
Physics Communications \textbf{188} (2015), 148--166.
\url{https://arxiv.org/abs/1403.3385}.

\bibitem{borcea}
J. Borcea and P. Br\"and\'en, \emph{The Lee--Yang and P\'olya--Schur
programs. II. Theory of stable polynomials and applications}.
\url{https://arxiv.org/abs/0809.3087}.
The positive-matrix stability argument needed here is proved in the text.

\bibitem{dlmf-hahn}
NIST Digital Library of Mathematical Functions, Chapter~18,
Sections~18.19--18.20, Hahn class and explicit representations.
\url{https://dlmf.nist.gov/18.20}. Accessed 11 October 2026 local time.

\bibitem{dlmf-zeta}
NIST Digital Library of Mathematical Functions, Section~25.11,
Hurwitz zeta function, shifts, derivatives, and Gamma relations.
\url{https://dlmf.nist.gov/25.11}. Accessed 11 October 2026 local time.

\end{thebibliography}
\end{document}
